%!TEX program = uplatex
\RequirePackage{plautopatch}
\documentclass[uplatex,dvipdfmx]{jlreq}
\usepackage{amsmath,amssymb}
\usepackage{url}

\pagestyle{empty}

\begin{document}

%======================================================================
% 表紙
%======================================================================
\begin{center}
  \vspace*{20pt}
  {\Huge\bfseries 双曲幾何}

  \vspace{1pc}
  {\Large 概要}
\end{center}

\vfill

\newpage

%======================================================================
% 講演１：小島定吉「双曲幾何入門」
%======================================================================
\noindent
{\large\bfseries 双曲幾何入門}
\hfill
{\large\bfseries 小島　定吉（東工大・情報理工）}

\bigskip

\noindent\textbf{要約：}

Bolyai、Lobachevsky に始まる双曲幾何は、
19世紀後半 Beltrami、Klein、Poincaré により
微分幾何学的解釈があたえられ数学の本流に現われ始めた。
この講演では、
現代数学の様々なスポットでの深遠な理論展開を語る
三氏の講演の前置きとして、
双曲幾何の基本的な事柄と素朴な感触を多少の歴史も込めて解説したい。

\vspace{6pt}

\begin{center}
  \textsc{参考文献}
\end{center}
\begin{enumerate}
  \item[{[1]}] 深谷賢治，「双曲幾何」，岩波講座 現代数学への入門, 岩波書店, 1996.
  \item[{[2]}] W.~サーストン（小島定吉 監訳），「3次元幾何学とトポロジー」，培風館, 1999.
  \item[{[3]}] 小島定吉，「多角形の現代幾何学 増補版」，牧野書店, 1999.
  \item[{[4]}] 小島定吉，「3次元の幾何学」，朝倉書店 講座 数学の考え方 22, 近刊.
  \item[{[5]}] J.~Stillwell,
        \textit{Sources of Hyperbolic Geometry},
        History of Mathematics \textbf{10},
        AMS-LMS, 1996.
\end{enumerate}

\newpage

%======================================================================
% 講演２：大鹿健一「双曲幾何と３次元トポロジー」
%======================================================================
\noindent
{\large\bfseries 双曲幾何と３次元トポロジー}
\hfill
{\large\bfseries 大鹿　健一（阪大・理）}

\bigskip

この話では、Thurston により導入された双曲幾何の手法が、
３次元位相幾何にもたらした結果を解説する。

３次元コンパクト多様体は素分解および Jaco--Shalen--Johannson により、
Seifert 多様体と呼ばれる、よく理解された空間と、atoroidal なものに分解される。
この atoroidal なものに双曲構造が入ることを示したのが Thurston の一意化定理である。

ここではこの定理の意味を述べることを第一の目標とし、
さらに重要な応用として Smith 予想の解決について解説する。

\newpage

%======================================================================
% 講演３：大鹿健一「双曲幾何と Klein 群」
%======================================================================
\noindent
{\large\bfseries 双曲幾何と Klein 群}
\hfill
{\large\bfseries 大鹿　健一（阪大・理）}

\bigskip

この話では、双曲幾何的手法が Klein 群の研究にもたらした成果を解説する。
まず Klein 群を双曲多様体の視点から捕らえることにより、
幾何的有限性などの概念が明確に定式化されるのを見る。

また永年懸案であった Ahlfors 予想のような難問が、双曲幾何をとおして、
位相幾何的問題に翻訳されることを解説する。
さらに双曲幾何的な Klein 群研究の発展の様子についても触れたい。

\newpage

%======================================================================
% 講演４：藤原耕二「双曲幾何と幾何学的群論」
%======================================================================
\noindent
{\large\bfseries 双曲幾何と幾何学的群論}
\hfill
{\large\bfseries 藤原　耕二（東北大・理）}

\bigskip

双曲幾何との関連を中心に、幾何学的群論について述べる。

\medskip
\noindent\textbf{I.}\quad
まず、Dehn の仕事について紹介する。
Dehn は前世紀始め、離散群（有限生成群）の3つの問題を定式化した。
すなわち、１）語の問題、２）共役問題、３）同型問題である。
Dehn はこれらを曲面群の場合に解決し、
その過程で quasi-isometry という概念も導入した。

\medskip
\noindent\textbf{II.}\quad
次に Gromov の双曲群について述べる。
定義と幾つかの例を挙げ、Dehn の３つの問題について述べる。
時間があれば、離散群の Dehn 関数、双曲群の外部自己同型群や
3次元多様体の双曲化との関連についても少し述べたい。

\medskip
\noindent\textbf{III.}\quad
最後に、有限表示群の JSJ 分解について述べる。
これは3次元多様体の Jaco--Shalen--Johannson 分解の導く基本群の分解を、
一般の有限表示群に拡張した理論である。
時間があれば、JSJ 理論の最近の応用についても述べたい。

\vspace{6pt}

\begin{center}
  \textsc{参考文献}
\end{center}

\noindent\textbf{I.}
\begin{enumerate}
  \item[{[LSc]}] R.~Lyndon and P.~Schupp,
        \textit{Combinatorial Group Theory},
        Springer-Verlag, 2001.
\end{enumerate}

\noindent\textbf{II.}
\begin{enumerate}
  \item[{[G]}] M.~Gromov,
        Hyperbolic groups,
        in \textit{Essays in Group Theory},
        Math.~Sci.~Res.~Inst.~Publ.~\textbf{8},
        Springer, 1987, 75--263.
  \item[{[S]}] Z.~Sela,
        The isomorphism problem for hyperbolic groups,
        \textit{Ann.~of Math.}~(2) \textbf{141} (1995), no.~2, 217--283.
  \item[{[O]}] 大鹿健一，離散群，岩波書店, 1998
               （岩波講座 現代数学の展開）.
  \item[{[F]}] 藤原耕二，「双曲群の手引き」，
               Surveys in Geometry 無限群と幾何学, 1995.\\
               \url{http://www.math.tohoku.ac.jp/~fujiwara/}
  \item[{[St]}] J.~R.~Stallings,
        On torsion-free groups with infinitely many ends,
        \textit{Ann.~of Math.}~(2) \textbf{88} (1968), 312--334.
\end{enumerate}

\noindent\textbf{III.}
\begin{enumerate}
  \item[{[RS]}] E.~Rips and Z.~Sela,
        Cyclic splittings of finitely presented groups and
        the canonical JSJ decomposition,
        \textit{Ann.~of Math.}~(2) \textbf{146} (1997), no.~1, 53--109.
  \item[{[Se]}] J.-P.~Serre,
        \textit{Trees},
        Springer-Verlag, 1980.
  \item[{[FP]}] K.~Fujiwara and P.~Papasoglu,
        JSJ-decompositions of finitely presented groups and complexes of groups,
        preprint.\\
        \url{http://www.math.tohoku.ac.jp/~fujiwara/}
  \item[{[S2]}] Z.~Sela,
        ICM 2002 Beijing, Proceedings, Algebra session.
\end{enumerate}

\newpage

%======================================================================
% 講演５：藤原一宏「数論と双曲幾何」
%======================================================================
\noindent
{\large\bfseries 数論と双曲幾何}
\hfill
{\large\bfseries 藤原　一宏（名古屋大・多元数理）}

\bigskip

双曲幾何学と数論は関係が深い。
しかしながら、このことはあまり多くの人々に認識されていないように見える
（一番よく知られているのは rigidity との絡みであるが、これについては述べない）。
この講演では、双方を結ぶ辞書を説明し、架け橋を造りたいと思っている。

テーマとしては

\medskip
\noindent\textbf{1.}\quad
非可換類体論の立場から見て三次元実双曲多様体は興味ある対象であり、
非可換相互法則を実現する場であること。
特に虚二次体の非可換相互法則と関連していること。

\medskip
\noindent\textbf{2.}\quad
Arithmetic topology との関連。
代数体の整数環や、有限体上の曲線はホモトピー的に見れば3次元であり、
双曲多様体の類似と思うのが自然である
（Mazur に始まる考えであり、最近の森下昌紀氏の仕事に続く）。

3次元幾何学における基本的な概念が数論における何に相当するか、
Wiles や肥田氏などが数論で行っているヘッケ環に関する仕事が
Thurston 理論にどう関係するかも述べたい。

\medskip

双方とも（特に2は）まだこれから発展することが期待される考え方であり、
分野にとらわれない見方をする利点を感じてもらいたい。

\vspace{6pt}

\begin{center}
  \textsc{参考文献}
\end{center}
\begin{enumerate}
  \item[{[1]}] 藤原一宏，
        論説「非可換類体論のめざすもの」，
        数学のたのしみ 17 号, 日本評論社.
  \item[{[2]}] 藤原一宏，
        論説「数学者の現実と夢」，
        数学のたのしみ 30 号, 日本評論社.
  \item[{[Hi]}] H.~Hida,
        On nearly ordinary Hecke algebras for $\mathrm{GL}(2)$
        over totally real fields,
        in \textit{Algebraic Number Theory},
        Adv.~Stud.~Pure Math.~\textbf{17},
        Academic Press, Boston, MA, 1989, 139--169.
  \item[{[Mo]}] M.~Morishita,
        Milnor's link invariants attached to certain Galois groups over $\mathbb{Q}$,
        \textit{Proc.~Japan Acad.~Ser.~A Math.~Sci.}~\textbf{76} (2000), no.~2, 18--21.
\end{enumerate}

\end{document}